\chapter{MATERIALS AND METHODS}

\section{Animal Preparation and Data Acquisition}
The empirical validation of the framework utilized the `c1608a01` dataset \cite{wunderle2013}. Multichannel electrophysiological recordings (Local Field Potentials and multi-unit spikes) were obtained from the primary visual cortex of an adult cat. To isolate intrinsic cortical network dynamics from confounding behavioral variables (e.g., saccadic eye movements or fluctuations in conscious attention), the animal was maintained under general anesthesia and paralysis. High-density matrix electrodes (32–64 channels) were implanted spanning Area 17 (A17), the A17/18 Transition Zone (TZ), and Area 18 (A18). 

\section{Data Ingestion Architecture (The VisionICeIO Pipeline)}
To process the high-density recordings, the \textit{VisionICeIO} pipeline \cite{visioniceio2024} (developed by Friedrich Schwarz, ORCID: \href{https://orcid.org/0009-0001-1167-8365}{0009-0001-1167-8365}) was integrated to act as a bridge between legacy proprietary formats (e.g., SPASS binaries) and modern data structures. 
\begin{itemize}
\item \textbf{Tensor Standardization:} Dispersed binaries containing spike times, waveforms, and continuous LFP traces were parsed and merged into regularized multidimensional arrays. Ragged arrays (e.g., varying spike counts per trial) were standardized using `NaN` padding.
\item \textbf{Zarr and Xarray Encapsulation:} The output is structured as a cloud-native \texttt{Zarr} store \cite{zarr2020}. By encapsulating the data within an \texttt{xarray.Dataset} \cite{hoyer2017xarray}, the framework leverages \textit{lazy loading}, allowing the frontend to read metadata trees without exhausting system RAM.
\end{itemize}

\section{Distributed Pre-Computation Engine (HPC)}
Analyzing dynamic connectivity across continuous recordings poses a severe combinatorial challenge. NECTA addresses this through a distributed pre-computation engine orchestrated by \textbf{Dask}\cite{rocklin2015dask}.
\begin{itemize}
\item \textbf{Task Graphs:} Computations are divided into discrete, independent tasks (e.g., computing coherence for a specific frequency band and time window) and mapped across a cluster of worker nodes.

\item \textbf{Fingerprinting and Cache Versioning:} To ensure data integrity, the orchestrator implements a cryptographic fingerprinting mechanism. Any modification to analytical parameters (e.g., altering variance thresholds or VAR criteria) alters the global hash, ensuring analytical consistency after parameter modifications.

\end{itemize}

\vspace{0.5cm}

% --- VISUAL INSERTION 2 USING TIKZ ---
\begin{figure}[htbp]
    \centering
    \begin{tikzpicture}[
        >=Stealth,
        block/.style={rectangle, draw, rounded corners, minimum width=3.2cm, minimum height=1.2cm, align=center, font=\sffamily\scriptsize\bfseries, thick},
        db/.style={cylinder, shape border rotate=90, draw, aspect=0.25, minimum width=2.5cm, minimum height=1.5cm, align=center, font=\sffamily\scriptsize\bfseries, thick, fill=green!10},
        worker/.style={rectangle, draw, rounded corners, minimum width=2.5cm, minimum height=1cm, align=center, fill=orange!10, font=\sffamily\scriptsize\bfseries, thick},
        arrow/.style={thick, ->, >=Stealth}
    ]

    % --- LAYER 1: INGESTION (XARRAY / ZARR) ---
    \node[block, fill=gray!10, dashed] (spass) at (-4.5, 0) {Legacy Binaries\\(SPASS LFP)};
    \node[block, fill=blue!10] (parser) at (0, 0) {VisionICeIO Pipeline\\(Data Ingestion)};
    \node[db] (zarr) at (4.5, 0.6) {Zarr Store\\(Data Chunks)};
    \node[block, fill=green!5] (xarray) at (4.5, -1.2) {Xarray Dataset\\(Metadata / Lazy)};

    \draw[arrow] (spass) -- node[above, font=\sffamily\tiny] {Parsing} (parser);
    \draw[arrow] (parser.east) -- (4.5, 0) -- (zarr.south);
    \draw[arrow] (parser.east) -- (4.5, 0) -- (xarray.north);
    \draw[<->, dashed, thick, color=green!50!black] (zarr.south) -- node[right, font=\sffamily\tiny] {Mapping} (xarray.north);

    % --- LAYER 2: DASK ORCHESTRATION ---
    \draw[draw=purple!80, thick, dashed, fill=purple!5, rounded corners] (-6, -3) rectangle (6, -8);
    \node[font=\sffamily\bfseries\scriptsize, purple!80!black, anchor=north] at (0, -3.1) {Dask Orchestration Infrastructure};

    \node[block, fill=white, draw=purple] (taskgraph) at (3.5, -4.5) {Task Graphs\\(Acyclic Graphs)};
    \node[block, fill=purple!10] (scheduler) at (-1.5, -4.5) {Dask Scheduler\\(Central Scheduler)};

    % Conexão Xarray -> Task Graph -> Scheduler
    \draw[arrow, ultra thick, color=green!60!black] (xarray.south) |- node[above right, font=\sffamily\tiny] {Trigger (Compute)} (taskgraph.north);
    \draw[arrow] (taskgraph.west) -- node[above, font=\sffamily\tiny] {Submission} (scheduler.east);

    % Workers
    \node[worker] (w1) at (-4, -6.5) {Worker Node 1\\(Chunk A)};
    \node[worker] (w2) at (0, -6.5) {Worker Node 2\\(Chunk B)};
    \node[worker] (w3) at (4, -6.5) {Worker Node N\\(Chunk N)};

    \draw[arrow] (scheduler.south) -- (-1.5, -5.5) -| (w1.north);
    \draw[arrow] (scheduler.south) -- (-1.5, -5.5) -| (w2.north);
    \draw[arrow] (scheduler.south) -- (-1.5, -5.5) -| (w3.north);

    % --- LAYER 3: AGGREGATION & FRONTEND ---
    \node[block, fill=orange!15, minimum width=6cm] (bus) at (0, -9.5) {Aggregation Bus\\(Results Synchronization)};
    \node[block, fill=red!10, draw=red!50!black, ultra thick, minimum width=6cm] (front) at (0, -11.5) {Reactive Frontend\\(Cytoscape / Dash)};

    % Routing from workers to bus
    \draw[arrow] (w1.south) |- (bus.west);
    \draw[arrow] (w2.south) -- (bus.north);
    \draw[arrow] (w3.south) |- (bus.east);
    
    \draw[arrow, ultra thick, draw=blue!70] (bus.south) -- node[right, font=\sffamily\scriptsize] {JSON Serialization} (front.north);

    \end{tikzpicture}
    
    \caption[Distributed and Cloud-Oriented Architecture]{\textbf{Block Diagram of Data Ingestion and Distribution in Dask.} The schematic details the architecture of the NECTA ecosystem: \textbf{(1)} Raw binaries are converted by the \textit{VisionICeIO Pipeline} into a chunked format, where heavy data resides in a \textit{Zarr Store} and the semantic structure (\textit{Metadata}) in an \textit{Xarray Dataset} that supports \textit{lazy loading}. \textbf{(2)} Upon requesting the computation, Xarray generates \textit{Task Graphs} (directed acyclic graphs) submitted to the \textit{Dask Scheduler}. \textbf{(3)} The scheduler stochastically distributes tasks (e.g., independent data \textit{Chunks}) across multiple \textit{Worker Nodes}. \textbf{(4)} The results computed in parallel are synchronized on a bus and exported via JSON to the interactive visualization layer of the \textit{Frontend}.}
    \label{fig:necta_architecture}
\end{figure}

\vspace{0.5cm}

\section{Spectral Connectivity and Phase Analysis}
To mitigate the artifacts of volume conduction inherent to LFP recordings, NECTA expands upon classical Magnitude Squared Coherence by employing phase-based estimators.
\begin{itemize}
\item \textbf{Analytical Signal Extraction:} Raw LFP signals are processed using the Hilbert Transform to yield the analytical signal, granting instantaneous access to both phase $\phi(t)$ and amplitude $A(t)$.

\item \textbf{Imaginary Coherence (iCoh) \& Weighted Phase Lag Index (wPLI):} The framework computes `iCoh` to strictly isolate interactions with non-zero phase delays. Furthermore, `wPLI` is calculated by weighting the phase differences by the magnitude of the imaginary component of the cross-spectrum, providing a highly robust metric against uncorrelated noise and instantaneous spatial spread.
\end{itemize}

\section{Directed Functional Connectivity (PDC)}
To infer causal interactions and hierarchical information flow, NECTA implements an Partial Directed Coherence (PDC) pipeline. Traditional PDC formulations rely on strictly causal MVAR models that forsake instantaneous effects, making them highly susceptible to the zero-lag correlations typical of volume conduction \cite{Faes_2010}. To overcome this, the NECTA framework utilizes:
\begin{enumerate}
    \item \textbf{Stationarity Filtering:} LFP windows exceeding a predefined variance threshold are identified as artifacts and discarded, preserving the assumptions of autoregressive modeling.
    
    \item \textbf{Frequency Transformation:} Autoregressive coefficients $A_r$ are transformed into the frequency domain $A(f)$ using Fast Fourier Transforms accelerated via \texttt{Numba} (FastMath) \cite{lam2015numba}, yielding the normalized PDC matrices.

    \item \textbf{Automated Order Selection:} Instead of imposing a static model order, the Dask workers dynamically fit Multivariate Autoregressive (MVAR) models using Ordinary Least Squares (OLS). The algorithm iterates through orders (up to \texttt{max\_order=20}) and selects the optimal $p$ that minimizes the Bayesian Information Criterion (BIC), balancing model fit against parametric complexity.
\end{enumerate}

Once connectivity matrices were estimated and statistically validated, graph-theoretical analyses were applied to characterize the resulting network topology.

\section{Graph Theory, Statistical Validation, and Open Science}
A crucial aspect of NECTA is its rigorous statistical and topological validation.

\begin{table}[hbtp]
\centering
\renewcommand{\arraystretch}{1.4} % Melhor espaçamento vertical para leitura
\footnotesize % Fonte menor para acomodar as 5 colunas de forma legível
\begin{tabular}{|>{\raggedright\arraybackslash}p{0.15\linewidth}|>{\raggedright\arraybackslash}p{0.16\linewidth}|>{\raggedright\arraybackslash}p{0.22\linewidth}|>{\raggedright\arraybackslash}p{0.18\linewidth}|>{\raggedright\arraybackslash}p{0.19\linewidth}|}
\hline
\textbf{Algorithmic Family} & \textbf{Representative Examples} & \textbf{Vulnerability to Instantaneous Spread (Volume Conduction)} & \textbf{Dimensionality of Causal Informational Direction} & \textbf{Justification Employed in NECTA Study} \\ \hline
\textit{Traditional Parametric Time-Domain Correlations} & Pearson Coefficient, Zero-lag covariance Cross-correlation & \textbf{High susceptibility.} Strong contamination risk due to the systematic generation of false instantaneous topological interactions. & Exclusively Symmetric and Non-Directed & Unfeasible for high-density mesoscopic analysis; method entirely suppressed from the Framework. \\ \hline
\textit{Mixed Spectral Synchrony} & Phase Locking Value (PLV), Classical Magnitude Squared Coherence & \textbf{Moderate susceptibility.} Latent topological interference in combined phase-magnitude domains. & Essentially Non-Directed & Baseline use restricted only for primary macro-anatomical recognition and historical comparison. \\ \hline
\textit{Analytical Modeling Delayed by Instantaneous Phase Shifts} & Weighted Phase Lag Index (wPLI), Imaginary Coherence (iCoh) & \textbf{Low susceptibility.} High robustness arising from consistent subtraction in complex domains of null real projection (consistent delay $>0$). & Non-Directed & Validated fundamental basis for clean extraction of network modularity sub-windows and artifact-free small-world coefficients. \\ \hline
\textit{Multi-Regressive Equation Models for Mathematical Causality} & Partial Directed Coherence (PDC), Directed Transfer Function (DTF) & \textbf{Medium to low susceptibility.} Attenuated and shielded in association and thresholding against non-null interactions via surrogate validation. & Exploratory Causal Asymmetry (Allows stipulating Driver/Source and Receiver/Sink) & Suitable for directional inference, coupled with optimized window processing via BIC, to track dynamic feline temporal hierarchical Feedforward/Feedback oscillation. \\ \hline
\end{tabular}
\caption{\textbf{Table of Methodological Synthesis of Connectivity Estimation.} Comparison of state-of-the-art metrics against the fundamental needs and physiological challenges faced by the NECTA framework.}
\label{tab:estimators}
\end{table}


\begin{itemize}
\item \textbf{Surrogate Data \& FDR Correction:} For connectivity edges, significance is established by generating 1000 phase-randomized surrogate datasets. Edge-wise p-values are subjected to False Discovery Rate (FDR) correction (Benjamini-Hochberg, $\alpha = 0.05$).

\item \textbf{Topological Null Models:} Macro-topology is validated against Erdös-Rényi and Configuration Models to ensure detected communities preserve empirical degree distributions.

\item \textbf{Open Science and Reproducibility:} Complying with open-source global standards, the NECTA framework facilitates transparency. All algorithms, Dask partition logs, and dependencies are containerized, version-controlled, and openly hosted on digital repositories (GitHub/CodeMeta), allowing independent researchers to audit, reproduce, and expand the methodological codebase.
\end{itemize}

\subsection{Topological Filtering: Maximum Spanning Tree and Proportional Thresholding}
\label{subsec:mst_thresholding}

The extraction of dense functional graphs (often resulting in fully connected networks) imposes a significant visual and analytical challenge, commonly known as the ``hairball effect'', which obscures the primary routes of information flow. To overcome this limitation and preserve the physiological integrity of the network, we implemented a dual topological filtering approach.

First, the Maximum Spanning Tree (MST) algorithm was applied using foundational graph theory libraries. The MST extracts the subgraph of maximum cumulative weight that connects all nodes in the system, ensuring that no recording channel (areas A17, A18, or the Transition Zone - TZ) becomes isolated from the connectome due to thresholding artifacts.

Following the structural anchoring via MST, the remaining edges are subjected to proportional thresholding (e.g., retaining the top 10\% to 15\% of the strongest connections). Mathematically, if $W$ is the coherence adjacency matrix, the filtered matrix $W'$ is defined as:

\begin{equation}
W'_{ij} = 
\begin{cases} 
W_{ij}, & \text{if } (i,j) \in \text{MST} \lor W_{ij} \geq P_{k} \\
0, & \text{otherwise}
\end{cases}
\end{equation}

where $P_{k}$ represents the $k$-th percentile of the edge weights. This combined approach effectively highlights the hub role of the Transition Zone (TZ) without compromising the global connectivity layout of the examined visual cortex.

\section{Interactive Visualization Dashboard}
The final tier of NECTA is a Single Page Application (SPA) built with \texttt{Dash} \cite{plotly2015} and \texttt{Cytoscape} \cite{shannon2003cytoscape}. The interface performs on-the-fly slicing of the pre-computed \texttt{Zarr} cache using \texttt{xarray.sel()}. Changes in the global state trigger instantaneous callbacks that serialize specific data slices into JSON, seamlessly re-rendering complex Connectome Matrices, Temporal Integration scatter plots, and anatomical Information Flow graphs.
